\section*{Chapitre \ref{ch:b2:mammal-reproduction} --- La reproduction sexuée des mammifères}

\begin{solution}{exo:b2:mammal-reproduction:1}
Les spermatogonies se divisent de la puberté à la fin de la vie~; les
ovogonies seulement chez le fœtus. Une \omterm{def:b2:meiosis-heredity:meiosis}{méiose} donne quatre
spermatozoïdes, mais un seul ovocyte (et trois globules polaires). Un
spermatozoïde demande environ 64 jours, plus deux semaines de
maturation~; un ovocyte reste bloqué en \omterm{prop:b2:meiosis-heredity:stages}{prophase I} d'avant la naissance
jusqu'au mois de son \omterm{def:b2:mammal-reproduction:ovary}{ovulation} --- des décennies --- puis de nouveau en
métaphase II jusqu'à la fécondation. Un homme produit $10^{8}$
spermatozoïdes par jour, quelque $10^{12}$ au cours d'une vie~; une
femme ovule environ 400 ovocytes.
\end{solution}

\begin{solution}{exo:b2:mammal-reproduction:2}
Fixation à la \omterm{prop:b2:mammal-reproduction:fertilisation}{zone pellucide}~; \omterm{prop:b2:mammal-reproduction:fertilisation}{réaction acrosomique} et digestion d'un
passage~; fusion des membranes et entrée du noyau du spermatozoïde~;
vague calcique et \omterm{prop:b2:mammal-reproduction:fertilisation}{réaction corticale} (durcissement de la zone, perte de
ses récepteurs~: le blocage de la polyspermie)~; achèvement de la
\omterm{def:b2:meiosis-heredity:meiosis}{méiose} II et expulsion du second globule polaire~; formation des deux
pronucléus, réplication de l'ADN, première mitose.
\end{solution}

\begin{solution}{exo:b2:mammal-reproduction:3}
Le \omterm{def:b2:mammal-reproduction:implantation}{trophoblaste} (couche externe) forme le \omterm{def:b2:mammal-reproduction:placenta}{placenta} et les annexes~; la
\omterm{def:b2:mammal-reproduction:implantation}{masse cellulaire interne} forme l'embryon. La \omterm{def:b2:mammal-reproduction:implantation}{nidation} commence le
6\textsuperscript{e} ou le 7\textsuperscript{e} jour après la
fécondation, dans une muqueuse utérine préparée par la progestérone du
\omterm{def:b2:mammal-reproduction:ovary}{corps jaune}.
\end{solution}

\begin{solution}{exo:b2:mammal-reproduction:4}
Dioxygène~: diffusion. Glucose~: transport facilité. Hématies
maternelles~: ne traversent pas (les sangs ne se mélangent jamais).
IgG~: transport par un récepteur. Urée~: diffusion, du fœtus vers la
mère. Alcool~: diffuse librement. Bactéries~: la plupart ne traversent
pas (certaines, comme celles de la syphilis et de la listériose, le
font).
\end{solution}

\begin{solution}{exo:b2:mammal-reproduction:5}
$m = 3, 1, 0.4, 0.1$~: $P = 0.95, 0.63, 0.33, 0.095$. Sous la moitié
lorsque $m < \ln 2 = 0.69$, c'est-à-dire en dessous de
$7\times 10^{7}$ spermatozoïdes. La courbe sature~: au-dessus d'environ
$10^{8}$, des spermatozoïdes supplémentaires n'ajoutent presque rien~; en
dessous, la fertilité diminue presque proportionnellement au nombre.
\end{solution}

\begin{solution}{exo:b2:mammal-reproduction:6}
Avant la puberté~: $7\times 10^{5}$ perdus en 13 ans (\num{4750}
jours)~: environ 150 par jour. Après~: \num{299000} perdus en 37 ans
(\num{13500} jours)~: 22 par jour~; \omterm{def:b2:mammal-reproduction:ovary}{ovulation} $400/\num{13500} = 0.03$
par jour. L'atrésie l'emporte sur l'\omterm{def:b2:mammal-reproduction:ovary}{ovulation} à sept cents contre un~;
le stock ne s'épuise pas par les \omterm{def:b2:mammal-reproduction:ovary}{ovulations} mais par la mort des
ovocytes.
\end{solution}

\begin{solution}{exo:b2:mammal-reproduction:7}
Adulte~: $(30/27)^{2.7} = 1.33$, $S = 1.33/2.33 = 0.57$. Fœtale~:
$(30/19)^{2.7} = 3.4$, $S = 3.4/4.4 = 0.77$. Aux basses pressions du
\omterm{def:b2:mammal-reproduction:placenta}{placenta}, le sang fœtal se charge d'un cinquième de dioxygène de plus
que ne le ferait le sang adulte, et le libère dans les tissus à des
pressions où l'hémoglobine adulte serait déjà à moitié vide.
\end{solution}

\begin{solution}{exo:b2:mammal-reproduction:8}
$2^{k} = 25$~: $k = 4.6$ doublements, \qty{9.3}{d}~: jour 16 à 17 après
la fécondation~; jour 30 à 31 après les dernières règles --- à peu près
le jour où les règles manquent.
\end{solution}

\begin{solution}{exo:b2:mammal-reproduction:9}
$3\times 10^{-8}\times 6\times 10^{4}\times 20/10^{-3} =
\qty{36}{mL/min}$ contre une demande de \qty{25}{mL/min}~: la marge
passe de huit à un et demi~; le fœtus devient hypoxique au moindre
effort supplémentaire, croît lentement et peut nécessiter un
accouchement prématuré.
\end{solution}

\begin{solution}{exo:b2:mammal-reproduction:10}
(a) Des \omterm{def:b2:mammal-reproduction:testis}{testicules} se forment et sécrètent les deux hormones~; l'AMH
élimine les canaux de Müller, mais la testostérone ne peut pas agir~:
les canaux de Wolff régressent et les organes génitaux externes sont
féminins --- une femme pourvue de \omterm{def:b2:mammal-reproduction:testis}{testicules} et sans utérus. (b) La
testostérone agit, l'AMH non~: un homme pourvu d'épididymes, de canaux
déférents et d'organes génitaux mâles, plus un utérus et des oviductes.
(c) Des ovaires et des dérivés müllériens (pas de \omterm{def:b2:mammal-reproduction:testis}{testicule}, pas
d'AMH)~; les androgènes surrénaliens virilisent plus ou moins les organes
génitaux externes. (d)
\emph{SRY} forme des \omterm{def:b2:mammal-reproduction:testis}{testicules}, qui imposent la voie mâle~: un homme à
deux chromosomes X, stérile car les gènes du Y nécessaires à la
\omterm{def:b2:mammal-reproduction:testis}{spermatogenèse} sont absents.
\end{solution}

\begin{solution}{exo:b2:mammal-reproduction:11}
Le cas le plus fréquent (\qty{68}{\%}) est la séparation entre les
jours 4 et 8~: après la formation du \omterm{def:b2:mammal-reproduction:implantation}{trophoblaste} (d'où un seul
\omterm{def:b2:mammal-reproduction:placenta}{placenta}), mais avant celle de l'amnios (d'où deux poches). Avant le
3\textsuperscript{e} jour, chaque moitié forme son propre \omterm{def:b2:mammal-reproduction:implantation}{trophoblaste},
d'où deux \omterm{def:b2:mammal-reproduction:placenta}{placentas} (\qty{30}{\%})~; après le 8\textsuperscript{e} jour,
l'amnios est déjà formé et il est partagé (\qty{2}{\%}). La division de
la \omterm{def:b2:mammal-reproduction:implantation}{masse cellulaire interne} après l'engagement du \omterm{def:b2:mammal-reproduction:implantation}{trophoblaste} est
manifestement l'événement le plus facile.
\end{solution}

\begin{solution}{exo:b2:mammal-reproduction:12}
Une femelle marsupiale investit peu avant la naissance~: une gestation
courte, un petit minuscule et une lactation qui peut être interrompue
en abandonnant le petit de la poche si la sécheresse survient~; son
risque par tentative est faible et elle peut aussitôt recommencer. Une
femelle placentaire investit lourdement dans une longue gestation d'où
naît un petit grand et bien développé, qui survit mieux à la naissance,
au prix d'une mère engagée pendant des mois et vulnérable entre-temps.
Dans les milieux productifs et prévisibles, la meilleure survie des
petits placentaires l'emporte~; le climat irrégulier de l'Australie
récompense la possibilité d'abandon des marsupiaux, et son isolement a
tenu les placentaires à l'écart jusqu'à une époque récente.
\end{solution}

\begin{solution}{pb:b2:mammal-reproduction:1}
\textbf{1.} $m = 3\times 10^{8}\times 10^{-8} = 3$~; $P = 1 - e^{-3} =
0.95$.
\textbf{2.} $1 - e^{-3}(1 + 3) = 0.80$~: sans le blocage, quatre
fécondations sur cinq seraient polyspermiques et l'embryon triploïde.
\textbf{3.} $0.15/5\times 10^{-5} = \qty{3000}{s}$, cinquante minutes~;
ce sont les contractions de l'utérus et de l'oviducte qui les
transportent.
\textbf{4.} $300/3\times 10^{8} = 10^{-6}$~: ils sont perdus dans le
vagin acide, dans la glaire cervicale, sous l'action des phagocytes de
l'utérus, et dans le mauvais oviducte.
\textbf{5.} De trois jours avant l'\omterm{def:b2:mammal-reproduction:ovary}{ovulation} à un jour après~: environ
quatre jours.
\textbf{6.} $m = 0.5$, $P = 0.39$~; nombre de cycles attendu $1/P = 2.5$.
\textbf{7.} Les quelques mitochondries du spermatozoïde, situées dans la
pièce intermédiaire, sont marquées et détruites après son entrée~;
l'ovocyte en porte cent mille.
\textbf{8.} $120/16 = 7.5$~: sept divisions, $2^{7} = 128$ cellules ---
environ une centaine.
\textbf{9.} $\log_2 25 = 4.6$ doublements, \qty{9.3}{d}~: jour 16 à 17~;
jour 30 à 31 après les règles.
\textbf{10.} Le \omterm{def:b2:mammal-reproduction:ovary}{corps jaune} régresse une douzaine de jours après
l'\omterm{def:b2:mammal-reproduction:ovary}{ovulation}, à moins que l'hCG ne le maintienne~; sans elle, la
progestérone chute, la muqueuse se desquame et l'embryon est perdu avec
les règles.
\textbf{11.} $\log_2 10^{5} = 16.6$ doublements~; vers la
10\textsuperscript{e} semaine, le \omterm{def:b2:mammal-reproduction:placenta}{placenta} produit sa propre
progestérone et le \omterm{def:b2:mammal-reproduction:ovary}{corps jaune} n'est plus nécessaire.
\textbf{12.} Les ovocytes restent en \omterm{prop:b2:meiosis-heredity:stages}{prophase I} pendant des décennies et
les cohésines qui maintiennent les bivalents se dégradent~: la
non-disjonction en \omterm{def:b2:meiosis-heredity:meiosis}{méiose} I augmente donc fortement avec l'âge
maternel~; les spermatozoïdes sont fabriqués en deux mois.
\textbf{13.} Après le \omterm{def:b2:mammal-reproduction:implantation}{trophoblaste} (jour 5) et avant l'amnios (jour
8)~: un \omterm{def:b2:mammal-reproduction:placenta}{placenta}, deux poches.
\textbf{14.} $3\times 10^{-8}\times 1.2\times 10^{5}\times 20/3.5\times
10^{-4} = \qty{206}{mL/min}$.
\textbf{15.} $7\times 3.5 = \qty{24.5}{mL/min}$~: une marge de huit.
\textbf{16.} $500\times 0.16 = \qty{80}{mL/min}$~; extraction
$24.5/80 = \qty{31}{\%}$.
\textbf{17.} $5\times 3.5\times 1440 = \qty{25}{g/d}$~; contenus dans
\qty{28}{L} de sang maternel~; débit journalier \qty{720}{L}~: le fœtus
prélève environ \qty{4}{\%} du glucose qui passe.
\textbf{18.} L'hémoglobine fœtale est saturée à \qty{77}{\%} à
\qty{30}{mmHg}, le fœtus a plus d'hématies et un débit cardiaque élevé,
et ses tissus sont adaptés à fonctionner à basse pression --- il vit, en
somme, au sommet de l'Everest.
\textbf{19.} Les échanges gazeux sans poumon fonctionnel, la nutrition
et l'excrétion sans intestin ni rein fonctionnels, et les hormones qui
maintiennent la grossesse~; il ne peut pas arrêter toutes les toxines
--- l'alcool, la nicotine et certains virus le traversent.
\textbf{20.} Étirement du col $\to$ hypothalamus $\to$ ocytocine $\to$
contraction $\to$ étirement accru, jusqu'à ce que l'expulsion supprime
le stimulus. Avant le terme, l'utérus dominé par la progestérone porte
peu de récepteurs à l'ocytocine et pas de jonctions communicantes~: les
contractions ne se propagent pas et la boucle ne peut pas se refermer.
\textbf{21.} $750\times 2.9 = \qty{2.2}{MJ}$~; coût $2.2/0.8 =
\qty{2.7}{MJ}$.
\textbf{22.} Croissance $30\times 6 = \qty{0.18}{MJ}$, entretien
$300\times 4 = \qty{1.2}{MJ}$~: \qty{1.4}{MJ}, environ les deux tiers de
l'énergie du lait~; le reste va à l'activité, à la thermorégulation et
aux pertes.
\textbf{23.} $270\times 2.7 = \qty{730}{MJ}$, deux fois et demie la
grossesse.
\textbf{24.} La tétée élève la prolactine et supprime la libération
pulsatile de l'hormone qui déclenche l'\omterm{def:b2:mammal-reproduction:ovary}{ovulation}~: l'\omterm{def:b2:mammal-reproduction:ovary}{ovulation} ne
revient que des mois plus tard~; sans contraception, cela espace les
naissances de deux à trois ans.
\textbf{25.} $P = 0.95$~; test positif le 17\textsuperscript{e} jour
après la fécondation (31\textsuperscript{e} jour du cycle)~; marge en
dioxygène de huit~; une journée de lait coûte \qty{2.7}{MJ} à la mère.
\end{solution}
